\chapter{Graphs: Strongly Connected Components, Bridges, and Articulation Points}
\chapterepigraph{This closing graph chapter asks about structural
weak points and tightly-knit clusters: which groups of nodes can all
reach each other (strongly connected components), and which single
edges or nodes -- if removed -- would break the graph apart (bridges
and articulation points)? Both questions are answered by one DFS,
carrying a little extra bookkeeping about \emph{when} each node was
discovered.}

\section{Discovery Time and Low-Link Value}
\label{sec:graph-others-intro}

Sections~\ref{sec:graph-bridges} and~\ref{sec:graph-articulation-points}
both rely on two numbers computed during a single DFS pass:

\begin{ideabox}[The Two Numbers Every Node Tracks]
\begin{itemize}
  \item $\textit{tin}[\textit{node}]$: the \textbf{discovery time} --
        the order in which DFS first visits this node (a simple
        counter, incremented on each new visit).
  \item $\textit{low}[\textit{node}]$: the \textbf{lowest discovery
        time reachable} from this node's DFS subtree, using any
        combination of tree edges (down to children) and \textbf{back
        edges} (up to an already-visited ancestor) -- but
        \textbf{never} the single tree edge back to the immediate
        parent, since that edge doesn't represent genuine alternate
        connectivity.
\end{itemize}
A node's $\textit{low}$ value answers: ``ignoring the edge I arrived
through, how far back up the DFS tree can I still reach?'' -- and that
single question, asked at every node, is what both bridges and
articulation points are built from.
\end{ideabox}

\begin{patternbox}[What This Chapter Covers]
\begin{itemize}
  \item \textbf{Strongly Connected Components (Section~\ref{sec:graph-kosaraju}):}
        a different technique entirely (Kosaraju's algorithm, using
        finish-time ordering and a graph transpose) for the directed-graph
        analogue of ``which nodes can all reach each other.''
  \item \textbf{Bridges in a Graph (Section~\ref{sec:graph-bridges})}:
        an edge $u\to v$ (tree edge) is a bridge exactly when
        $v$'s subtree has \textbf{no} back edge reaching $u$ or higher
        -- i.e.\ $\textit{low}[v] > \textit{tin}[u]$.
  \item \textbf{Articulation Points (Section~\ref{sec:graph-articulation-points})}:
        a near-identical condition, $\textit{low}[v] \ge
        \textit{tin}[u]$ (note: $\ge$, not $>$), identifying
        \textbf{nodes} whose removal disconnects the graph, with one
        special case for the DFS root.
\end{itemize}
\end{patternbox}

\newpage

\section{Strongly Connected Components (Kosaraju's Algorithm)}
\label{sec:graph-kosaraju}

\problemstatement{Given a directed graph, find its \textbf{Strongly
Connected Components} (SCCs) -- maximal groups of nodes where every
node can reach every other node in the same group via directed
edges.}

\begin{patternbox}
\emph{``Group nodes that can all reach each other''} is solved by
Kosaraju's algorithm: run DFS once to compute a finish-time order
(exactly the Graph Topological Sort chapter's finish-time idea), then
run DFS \textbf{again} on the graph with every edge \textbf{reversed},
processing nodes in \textbf{decreasing} finish-time order -- each DFS
launched in this second pass visits exactly one complete SCC.
\end{patternbox}

\subsection*{Why reversing edges and using finish-time order works}

Reversing every edge preserves SCCs exactly (if $u$ and $v$ can reach
each other in the original graph, they can still reach each other with
every edge flipped) but \textbf{destroys} inter-SCC reachability in a
specific, useful direction: collapsing each SCC into a single
``super-node'' produces a DAG (the \textbf{condensation} of the
graph), and processing nodes in decreasing finish-time order from the
\emph{first} DFS guarantees the first unvisited node encountered in
the second pass always belongs to a \textbf{source} SCC (no incoming
edges from any other not-yet-visited SCC) of the \emph{original}
condensation -- which is a \textbf{sink} SCC once edges are reversed.
Starting DFS there in the reversed graph cannot leak into any other
SCC, since reversed edges now only point \emph{into} this SCC from
elsewhere, not out of it.

\subsection*{Algorithm}

\begin{enumerate}
  \item Run DFS over the original graph from every unvisited node,
        pushing each node onto a stack when it finishes (identical to
        topological sort's DFS approach).
  \item Build the transposed graph (every edge reversed).
  \item Pop the stack one node at a time; for every still-unvisited
        node, run DFS on the \textbf{transposed} graph -- everything
        reached in this one DFS call is exactly one SCC.
\end{enumerate}

\subsection*{C++ implementation}

\begin{lstlisting}
#include <bits/stdc++.h>
using namespace std;

void fillOrder(int node, vector<vector<int>>& adj, vector<bool>& visited, stack<int>& order) {
    visited[node] = true;
    for (int neighbor : adj[node]) {
        if (!visited[neighbor]) fillOrder(neighbor, adj, visited, order);
    }
    order.push(node);
}

void collectSCC(int node, vector<vector<int>>& transposed, vector<bool>& visited,
                 vector<int>& component) {
    visited[node] = true;
    component.push_back(node);
    for (int neighbor : transposed[node]) {
        if (!visited[neighbor]) collectSCC(neighbor, transposed, visited, component);
    }
}

vector<vector<int>> kosaraju(int n, vector<vector<int>>& adj) {
    vector<bool> visited(n, false);
    stack<int> order;
    for (int i = 0; i < n; i++) {
        if (!visited[i]) fillOrder(i, adj, visited, order);
    }

    vector<vector<int>> transposed(n);
    for (int u = 0; u < n; u++) {
        for (int v : adj[u]) transposed[v].push_back(u);
    }

    fill(visited.begin(), visited.end(), false);
    vector<vector<int>> sccs;
    while (!order.empty()) {
        int node = order.top();
        order.pop();
        if (!visited[node]) {
            vector<int> component;
            collectSCC(node, transposed, visited, component);
            sccs.push_back(component);
        }
    }
    return sccs;
}

int main() {
    int n = 5;
    vector<vector<int>> adj(n);
    adj[0] = {2}; adj[2] = {1}; adj[1] = {0}; adj[0].push_back(3); adj[3] = {4};

    for (auto& scc : kosaraju(n, adj)) {
        for (int x : scc) cout << x << " ";
        cout << endl;
    }
    // 0 1 2   (one SCC)
    // 3        (its own SCC)
    // 4        (its own SCC)
    return 0;
}
\end{lstlisting}

\subsection*{Dry run}

\dryrun{Take edges $0{\to}2,2{\to}1,1{\to}0,0{\to}3,3{\to}4$.}

First DFS (finish-time order, starting from $0$): visits $0\to2\to1$
(back to $0$, already visiting, skip), $1$ finishes first (push),
$2$ finishes (push), then $0$ continues to $3\to4$, $4$ finishes
(push), $3$ finishes (push), $0$ finishes (push). Stack
(bottom-to-top): $1,2,4,3,0$. Transposed graph: $2{\to}0,1{\to}2,0{\to}1,
3{\to}0,4{\to}3$. Second DFS, popping stack top-to-bottom: pop $0$
first -- DFS on transposed graph from $0$: $0\to1\to2\to0$ (visited,
stop) -- component $\{0,1,2\}$. Pop $3$ (next unvisited): DFS from $3$
on transposed graph: $3\to0$ (visited, stop) -- component $\{3\}$.
Pop $4$: component $\{4\}$. Three SCCs: $\{0,1,2\}$, $\{3\}$, $\{4\}$
-- correctly identifying the $3$-cycle as one SCC and the two
remaining nodes (no return path) as singleton SCCs.

\begin{complexitybox}
\textbf{Time Complexity.} $O(V+E)$ -- two full DFS passes plus
building the transpose, each linear in the graph size.
\textbf{Space Complexity.} $O(V+E)$ for the transposed adjacency list,
plus $O(V)$ for the stack and visited arrays.
\end{complexitybox}

\begin{takeawaybox}
Kosaraju's algorithm is this book's clearest demonstration that
\textbf{reversing} a directed graph's edges is sometimes itself the
key algorithmic move -- not just bookkeeping, but the step that turns
an otherwise hard global structure question into two simple, linear
DFS passes.
\end{takeawaybox}

\section{Bridges in a Graph}
\label{sec:graph-bridges}

\problemstatement{Given an undirected, connected graph, find every
\textbf{bridge} -- an edge whose removal would increase the number of
connected components (disconnect the graph).}

\begin{patternbox}
\emph{``Find edges whose removal disconnects the graph''} is solved by
a single DFS tracking $\textit{tin}$ and $\textit{low}$
(Section~\ref{sec:graph-others-intro}): a tree edge $u \to v$ (where
$v$ is $u$'s child in the DFS tree) is a bridge exactly when
$\textit{low}[v] > \textit{tin}[u]$ -- meaning $v$'s entire subtree
has \textbf{no} back edge reaching $u$ or any of $u$'s ancestors, so
the only way to reach $v$'s subtree from the rest of the graph is
through this one edge.
\end{patternbox}

\subsection*{Deriving the \textit{low}-value update rule}

While DFS explores from node $u$, for each neighbor $v$:
\begin{itemize}
  \item If $v$ is $u$'s \textbf{parent} (the edge just arrived
        through), skip it entirely -- this is not a genuine
        alternative path.
  \item If $v$ is \textbf{unvisited}, it becomes a tree-edge child:
        recurse, then update $\textit{low}[u] \gets
        \min(\textit{low}[u], \textit{low}[v])$ (whatever $v$'s
        subtree can reach back to, $u$ can also reach back to, through
        $v$). Check the bridge condition here:
        $\textit{low}[v] > \textit{tin}[u]$.
  \item If $v$ is \textbf{already visited} (and not the parent), this
        is a back edge: update $\textit{low}[u] \gets
        \min(\textit{low}[u], \textit{tin}[v])$ -- using $v$'s
        \textbf{discovery time} directly (not $\textit{low}[v]$, since
        crossing a back edge to $v$ only guarantees reaching $v$
        itself, not necessarily anything $v$'s own low value might
        represent through edges not yet explored from $u$'s
        perspective).
\end{itemize}

\subsection*{C++ implementation}

\begin{lstlisting}
#include <bits/stdc++.h>
using namespace std;

int timer = 0;

void dfs(int node, int parent, vector<vector<int>>& adj, vector<bool>& visited,
         vector<int>& tin, vector<int>& low, vector<vector<int>>& bridges) {
    visited[node] = true;
    tin[node] = low[node] = timer++;

    for (int neighbor : adj[node]) {
        if (neighbor == parent) continue;

        if (visited[neighbor]) {
            low[node] = min(low[node], tin[neighbor]); // back edge
        } else {
            dfs(neighbor, node, adj, visited, tin, low, bridges);
            low[node] = min(low[node], low[neighbor]);
            if (low[neighbor] > tin[node]) {
                bridges.push_back({node, neighbor});
            }
        }
    }
}

vector<vector<int>> findBridges(int n, vector<vector<int>>& adj) {
    vector<bool> visited(n, false);
    vector<int> tin(n), low(n);
    vector<vector<int>> bridges;
    timer = 0;

    for (int i = 0; i < n; i++) {
        if (!visited[i]) dfs(i, -1, adj, visited, tin, low, bridges);
    }
    return bridges;
}

int main() {
    int n = 4;
    vector<vector<int>> adj(n);
    auto addEdge = [&](int u, int v) { adj[u].push_back(v); adj[v].push_back(u); };
    addEdge(0, 1); addEdge(1, 2); addEdge(2, 0); addEdge(1, 3);

    for (auto& b : findBridges(n, adj)) cout << b[0] << "-" << b[1] << " ";
    cout << endl; // 1-3
    return 0;
}
\end{lstlisting}

\subsection*{Dry run}

\dryrun{Take edges $0{-}1,1{-}2,2{-}0,1{-}3$ (a triangle $0,1,2$ with
a pendant edge to $3$), DFS from $0$.}

$\textit{dfs}(0,-1)$: $\textit{tin}[0]=\textit{low}[0]=0$. Neighbor
$1$: unvisited, recurse. $\textit{dfs}(1,0)$:
$\textit{tin}[1]=\textit{low}[1]=1$. Neighbor $0$: is parent, skip.
Neighbor $2$: unvisited, recurse. $\textit{dfs}(2,1)$:
$\textit{tin}[2]=\textit{low}[2]=2$. Neighbor $1$: is parent, skip.
Neighbor $0$: already visited, not parent -- back edge,
$\textit{low}[2]=\min(2,\textit{tin}[0]{=}0)=0$. Return to $1$:
$\textit{low}[1]=\min(1,\textit{low}[2]{=}0)=0$; check bridge:
$\textit{low}[2]{=}0 > \textit{tin}[1]{=}1$? No -- not a bridge
(correct: edge $1{-}2$ sits on the triangle's cycle). Back at $1$,
next neighbor $3$: unvisited, recurse. $\textit{dfs}(3,1)$:
$\textit{tin}[3]=\textit{low}[3]=3$; no unvisited or back-edge
neighbors. Return to $1$: $\textit{low}[1]=\min(0,\textit{low}[3]{=}3)=0$;
check bridge: $\textit{low}[3]{=}3 > \textit{tin}[1]{=}1$? \textbf{Yes}
-- edge $1{-}3$ is a bridge (correct: it's the only connection to the
otherwise-isolated node $3$). Return to $0$: $\textit{low}[0]=
\min(0,\textit{low}[1]{=}0)=0$; check bridge:
$\textit{low}[1]{=}0 > \textit{tin}[0]{=}0$? No -- not a bridge.

\begin{complexitybox}
\textbf{Time Complexity.} $O(V+E)$ -- a single DFS pass.
\textbf{Space Complexity.} $O(V)$ for \textit{tin}, \textit{low},
\textit{visited}, and the recursion stack.
\end{complexitybox}

\begin{takeawaybox}
The distinction between updating $\textit{low}$ with a child's
$\textit{low}$ value (tree edges) versus a neighbor's $\textit{tin}$
value (back edges) is the single most important detail to get right in
this algorithm -- using $\textit{low}$ for a back edge would
incorrectly let information ``leak'' through a node that hasn't
finished exploring its own subtree yet.
\end{takeawaybox}

\section{Articulation Points}
\label{sec:graph-articulation-points}

\problemstatement{Given an undirected, connected graph, find every
\textbf{articulation point} (cut vertex) -- a node whose removal
(along with all its incident edges) would increase the number of
connected components.}

\begin{patternbox}
This closing problem of the entire graph series reuses
Section~\ref{sec:graph-bridges}'s exact $\textit{tin}$/$\textit{low}$
DFS, with the condition shifted from edges to \textbf{nodes}, and from
strict to non-strict inequality: a non-root node $u$ is an
articulation point if it has a child $v$ with
$\textit{low}[v] \ge \textit{tin}[u]$ -- plus one special case for the
DFS \textbf{root}, which has no parent for the usual argument to apply
to.
\end{patternbox}

\subsection*{Why $\ge$ replaces bridges' strict $>$}

A bridge requires $v$'s subtree to have \textbf{no} way back to $u$ or
higher ($\textit{low}[v] > \textit{tin}[u]$, strictly). An articulation
point is a weaker condition: $u$ is critical even if $v$'s subtree can
reach back \emph{exactly} to $u$ itself ($\textit{low}[v] =
\textit{tin}[u]$) -- because reaching back only as far as $u$ still
means removing $u$ disconnects $v$'s subtree from everything
\emph{above} $u$, even though $v$'s subtree remains internally
connected to $u$'s position. Only if $v$'s subtree can reach back
\emph{above} $u$ (to some ancestor, $\textit{low}[v] < \textit{tin}[u]$)
does removing $u$ fail to disconnect anything.

\subsection*{The root's special case}

The root of the DFS tree has no parent, so the ``can my child's
subtree reach above me'' question doesn't apply to it the same way.
Instead: the root is an articulation point exactly when it has
\textbf{more than one child} in the DFS tree -- each child starts an
independent branch that (by definition of being a separate DFS tree
child rather than reached via a back edge) has no other connection to
the rest of the graph except through the root.

\subsection*{C++ implementation}

\begin{lstlisting}
#include <bits/stdc++.h>
using namespace std;

int timer = 0;

void dfs(int node, int parent, vector<vector<int>>& adj, vector<bool>& visited,
         vector<int>& tin, vector<int>& low, vector<bool>& isArticulation) {
    visited[node] = true;
    tin[node] = low[node] = timer++;
    int children = 0;

    for (int neighbor : adj[node]) {
        if (neighbor == parent) continue;

        if (visited[neighbor]) {
            low[node] = min(low[node], tin[neighbor]);
        } else {
            dfs(neighbor, node, adj, visited, tin, low, isArticulation);
            low[node] = min(low[node], low[neighbor]);

            if (parent != -1 && low[neighbor] >= tin[node]) {
                isArticulation[node] = true;
            }
            children++;
        }
    }

    if (parent == -1 && children > 1) {
        isArticulation[node] = true;
    }
}

vector<int> findArticulationPoints(int n, vector<vector<int>>& adj) {
    vector<bool> visited(n, false), isArticulation(n, false);
    vector<int> tin(n), low(n);
    timer = 0;

    for (int i = 0; i < n; i++) {
        if (!visited[i]) dfs(i, -1, adj, visited, tin, low, isArticulation);
    }

    vector<int> result;
    for (int i = 0; i < n; i++) {
        if (isArticulation[i]) result.push_back(i);
    }
    return result;
}

int main() {
    int n = 5;
    vector<vector<int>> adj(n);
    auto addEdge = [&](int u, int v) { adj[u].push_back(v); adj[v].push_back(u); };
    addEdge(0, 1); addEdge(1, 2); addEdge(2, 0); addEdge(1, 3); addEdge(3, 4);

    for (int x : findArticulationPoints(n, adj)) cout << x << " ";
    cout << endl; // 1 3
    return 0;
}
\end{lstlisting}

\subsection*{Dry run}

\dryrun{Take edges $0{-}1,1{-}2,2{-}0,1{-}3,3{-}4$ (a triangle
$0,1,2$, with $1$ also connected via $3$ to a pendant $4$), DFS from
$0$.}

This is the same graph shape as Section~\ref{sec:graph-bridges}'s
example with an added pendant ($3{-}4$). Following the same DFS order:
$\textit{tin}/\textit{low}$ for $0,1,2$ resolve identically (the
triangle is unaffected). At node $3$ (child of $1$, with one child of
its own, $4$): $\textit{dfs}(4,3)$ gives
$\textit{tin}[4]=\textit{low}[4]=4$. Back at $3$:
$\textit{low}[3]=\min(3,4)=3$; check articulation: parent of $3$ is
$1$ ($\neq-1$), $\textit{low}[4]{=}4 \ge \textit{tin}[3]{=}3$? Yes --
\textbf{$3$ is an articulation point} (removing it disconnects $4$).
Back at $1$ (processing child $3$): $\textit{low}[1]=\min(0,
\textit{low}[3]{=}3)=0$; check articulation: parent of $1$ is $0$
($\neq-1$), $\textit{low}[3]{=}3 \ge \textit{tin}[1]{=}1$? Yes --
\textbf{$1$ is an articulation point} (removing it disconnects $\{3,4\}$
from the triangle). Node $0$ (the root): only one DFS-tree child
($1$; node $2$ was reached via a back edge, not as a tree child), so
the root special case does \textbf{not} apply -- $0$ is not an
articulation point. Final: $\{1,3\}$.

\begin{complexitybox}
\textbf{Time Complexity.} $O(V+E)$ -- a single DFS pass, identical
structure to bridge-finding.
\textbf{Space Complexity.} $O(V)$ for \textit{tin}, \textit{low},
\textit{visited}, \textit{isArticulation}, and the recursion stack.
\end{complexitybox}

\begin{takeawaybox}
This book series closes its graph coverage on a fitting note: bridges
and articulation points are, almost letter for letter, the same
algorithm, distinguished only by a strict-versus-non-strict inequality
and one root-specific special case. Across every chapter in this
series -- Greedy, Binary Search, Heap, Stack/Queue, BST, Trie,
Recursion, eight DP sub-topics, and six graph sub-topics -- the
recurring lesson has been the same: a small number of genuinely new
ideas, recognized and reapplied with small, deliberate
modifications, cover an enormous space of distinct-looking interview
problems.
\end{takeawaybox}

